\section{Introduction}\label{sec:intro}
% Budget: ~1.25 pp (abstract + intro <= 1.5 pp). Fig.~1 lives here.
% Checklist (facts to cover, author writes the text):
%  - Setting: logged bandit data accumulated over many periods; want V_T(pi) for the
%    current/target period; only a small anchor sample from T.
%  - Tension: all history = low variance, biased under drift; recent/anchor only = unbiased,
%    high variance. Standard fixes (recency windows, decay, forecasting) treat all drift alike.
%  - Key observation: context drift and logger drift are correctable (transport /
%    propensities); only conditional-reward drift q_t(x,a) biases a reweighted estimate.
%  - Consequence: test for reward drift on the conditional reward (per-action residual test),
%    not on the scalar policy value (too noisy; PLAN item 6-9), and reuse old segments that
%    pass, with a correction for those that do not.
%  - Contributions (3 bullets max): (i) mechanism decomposition + detection on q;
%    (ii) estimator: segment-wise joint reward models, reference selection by anchor,
%    drift-corrected reuse, min-variance pooling, random-effects interval;
%    (iii) pre-registered evaluation: simulator regimes S0--S8 + Open Bandit
%    semi-synthetic, one-shot test split, reported failure cases.
%  - Honest scope sentence: where it does not help (S3b gradual drift; misspecified
%    reward drift in S4 -- pending v3 decision).
\todo{introduction}

% Fig. 1: mechanism illustration -- context-shifted old segment kept, reward-shifted
% segment corrected/discounted. Source: figures/fig1_mechanism.pdf (to make).
